\begin{figure}[H]
\centering
\includegraphics[width=1.0\linewidth]{wakayama/F7-1_wind_roses.pdf}
\caption[風配図（前半）]{代表格子点の風配図、\textbf{月の前半}（1--15 日）。
平均と楕円では見せられない「分布そのもの」を示す。
ローブが 2 つに割れている場面が、平均風がほぼゼロで分散だけが大きくなる時期にあたる。}
\label{fig:f7-wakayama-h1}
\end{figure}

\begin{figure}[H]
\centering
\includegraphics[width=1.0\linewidth]{wakayama/F7-2_wind_roses.pdf}
\caption[風配図（後半）]{同じ格子点の\textbf{月の後半}（16 日--月末）。
前半（図\ref{fig:f7-wakayama-h1}）と並べると、ほとんどの月で\textbf{2 枚はよく似ている}。
唯一 10 月だけ \SI{30}{\hecto\pascal} の最頻セクターが前半 \ang{90}、
後半 \ang{248} と入れ替わるが、どちらも頻度は 2 割前後しかない。
つまり向きが変わったというより、\textbf{その時期はどちらの半月も
風向が定まっていない}のであって、最頻セクターを読む意味が薄い。}
\label{fig:f7-wakayama-h2}
\end{figure}